% !TEX program = xelatex
\documentclass[14pt,a4paper]{extarticle}

\usepackage{fontspec}
\usepackage{polyglossia}
\setdefaultlanguage{russian}
\setotherlanguage{english}
\IfFontExistsTF{Times New Roman}{\setmainfont{Times New Roman}}{\IfFontExistsTF{Liberation Serif}{\setmainfont{Liberation Serif}}{\setmainfont{FreeSerif}}}
\IfFontExistsTF{Liberation Mono}{\setmonofont{Liberation Mono}}{\setmonofont{FreeMono}}
\IfFontExistsTF{Times New Roman}{\newfontfamily\cyrillicfont{Times New Roman}[Script=Cyrillic]}{\newfontfamily\cyrillicfont{Liberation Serif}[Script=Cyrillic]}

\usepackage[a4paper,left=30mm,right=15mm,top=20mm,bottom=20mm]{geometry}
\usepackage{setspace}
\usepackage{indentfirst}
\usepackage{amsmath,amssymb}
\usepackage{graphicx}
\usepackage{float}
\usepackage{array}
\usepackage{longtable}
\usepackage{booktabs}
\usepackage{caption}
\usepackage{hyperref}
\usepackage{xurl}
\usepackage{enumitem}
\usepackage{makecell}
\usepackage{multirow}

\hypersetup{hidelinks}
\onehalfspacing
\setlength{\parindent}{1.25cm}
\setlength{\parskip}{0pt}
\captionsetup{labelsep=endash, justification=centering}
\renewcommand{\arraystretch}{1.15}
\sloppy

\newcommand{\figHistPath}{assets/histogram_receipts.png}
\newcommand{\figLognormPath}{assets/histogram_lognormal_fit.png}
\newcommand{\figTopPath}{assets/top10_receipts.png}
\newcommand{\figComparePath}{assets/change_2019_2021.png}
\newcommand{\figBoxPath}{assets/boxplot_receipts.png}
\newcommand{\insertfig}[2][0.88\textwidth]{%
    \IfFileExists{#2}{\includegraphics[width=#1]{#2}}{%
        \fbox{\parbox[c][7cm][c]{0.82\textwidth}{\centering
        Файл изображения не найден.\\
        Проверьте путь: \texttt{\detokenize{#2}}}}%
    }%
}

\begin{document}

\tableofcontents
\newpage

\section*{Введение}
\addcontentsline{toc}{section}{Введение}

В работе рассматривается показатель \textit{Receipts per Arrival}, то есть средний объем туристских поступлений, приходящийся на одно международное прибытие. Этот показатель удобен тем, что позволяет оценить не только масштаб туристского потока, но и средний объем поступлений, который приходится на одно прибытие. Иначе говоря, он помогает отделить массовость направления от доходности одной поездки.

Для исследования выбран 2021 год. Этот период интересен тем, что туристский рынок уже начал восстанавливаться после наиболее жестких ограничений, однако структура поездок еще оставалась неоднородной. Поэтому можно ожидать заметный разброс между странами как по числу прибытий, так и по среднему объему расходов на одну поездку.

Цель работы - провести статистический анализ показателя \textit{Receipts per Arrival} по странам Европы, построить наглядные графики, определить основные числовые характеристики, проверить гипотезу о законе распределения и сформулировать содержательные выводы. Кроме того, в работе выполнено сравнение стран между собой, проведена укрупненная группировка стран и выполнено сопоставление с 2019 годом.

\section{Основная часть}

\subsection{Самостоятельный поиск статистических данных}

Для выполнения практической работы были использованы материалы международной базы UN Tourism. Поиск данных выполнялся самостоятельно на официальном сайте организации в разделе \textit{Tourism Statistics Database - Inbound Tourism}. Внутри базы были выбраны два массива данных: по международным туристским прибытиям и по поступлениям от въездного туризма.

Готового показателя \textit{Receipts per Arrival} в исходных таблицах нет, поэтому он был рассчитан самостоятельно. Для этого использованы:
\begin{enumerate}[leftmargin=1.25cm]
    \item таблица \textit{Inbound Tourism: Arrivals};
    \item таблица \textit{Inbound Tourism: Expenditure}.
\end{enumerate}

В итоговую выборку включались только те страны, для которых одновременно имелись оба нужных значения за 2021 год: число прибытий и общий объем поступлений. После отбора была получена выборка объемом $n=41$ наблюдение. В нее вошли страны европейского туристского пространства, представленные в базе UN Tourism и имеющие сопоставимые данные за рассматриваемый год.

\subsection{Описание используемых данных}

В таблице по прибытиям использован показатель \textit{inbound -- trips -- total -- total -- overnight visitors (tourists)}. Он показывает количество международных прибытий туристов с ночевкой. Значения представлены в тысячах прибытий.

В таблице по поступлениям использован показатель \textit{inbound -- expenditure -- balance of payments -- total -- visitors}. Он отражает суммарные поступления от въездного туризма. Значения приведены в миллионах долларов США.

Искомый показатель для $i$-й страны рассчитывался по формуле
\begin{equation}
R_i = \frac{E_i}{A_i} \cdot 1000,
\end{equation}
где $R_i$ - \textit{Receipts per Arrival}, USD на одно прибытие; $E_i$ - поступления от въездного туризма, млн USD; $A_i$ - число прибытий, тыс. человек.

Множитель $1000$ появляется из-за того, что в исходных таблицах используются разные единицы измерения: миллионы долларов и тысячи прибытий. В результате после пересчета получается показатель в долларах США на одно прибытие.

Основные источники данных:
\begin{enumerate}[leftmargin=1.25cm]
    \item сайт UN Tourism: \url{https://www.untourism.int/tourism-statistics/tourism-data-inbound-tourism};
    \item таблица по прибытиям: \url{https://pre-webunwto.s3.amazonaws.com/s3fs-public/2025-12/UN_Tourism_inbound_arrivals_12_2025.xlsx};
    \item таблица по поступлениям: \url{https://pre-webunwto.s3.amazonaws.com/s3fs-public/2025-12/UN_Tourism_inbound_expenditure_12_2025.xlsx}.
\end{enumerate}

\subsection{Формальное представление исследуемых данных}

Подготовленная выборка приведена в таблице~\ref{tab:data}. В ней для каждой страны показаны исходные величины и рассчитанный итоговый показатель. Таблица приведена полностью, поскольку объем выборки не слишком велик, а полный перечень стран нужен для дальнейшей интерпретации результатов.

\small
\begin{longtable}{|c|p{4.4cm}|>{\raggedleft\arraybackslash}p{2.2cm}|>{\raggedleft\arraybackslash}p{2.4cm}|>{\raggedleft\arraybackslash}p{2.5cm}|}
\caption{Показатель \textit{Receipts per Arrival} по странам Европы за 2021 год} \label{tab:data}\\
\hline
№ & Страна & Прибытия, тыс. & Поступления, млн USD & Receipts per Arrival, USD \\ \hline
\endfirsthead
\hline
№ & Страна & Прибытия, тыс. & Поступления, млн USD & Receipts per Arrival, USD \\ \hline
\endhead
1 & Австрия & 12728,00 & 11944,47 & 938,44 \\ \hline
2 & Албания & 5515,00 & 2479,95 & 449,67 \\ \hline
3 & Андорра & 1949,00 & 1924,88 & 987,62 \\ \hline
4 & Армения & 876,00 & 813,43 & 928,57 \\ \hline
5 & Беларусь & 787,00 & 661,53 & 840,57 \\ \hline
6 & Бельгия & 3243,00 & 8305,29 & 2560,99 \\ \hline
7 & Болгария & 2300,00 & 2720,34 & 1182,76 \\ \hline
8 & Босния и Герцеговина & 502,00 & 1061,22 & 2113,98 \\ \hline
9 & Великобритания & 6287,00 & 37790,49 & 6010,89 \\ \hline
10 & Венгрия & 7929,10 & 5692,55 & 717,93 \\ \hline
11 & Германия & 11652,00 & 22137,26 & 1899,87 \\ \hline
12 & Греция & 14705,00 & 13666,11 & 929,35 \\ \hline
13 & Грузия & 1577,00 & 1350,61 & 856,44 \\ \hline
14 & Дания & 7555,00 & 4426,76 & 585,94 \\ \hline
15 & Исландия & 698,00 & 1302,80 & 1866,48 \\ \hline
16 & Испания & 31181,00 & 34176,16 & 1096,06 \\ \hline
17 & Италия & 26888,00 & 25354,93 & 942,98 \\ \hline
18 & Кипр & 1936,93 & 2099,71 & 1084,04 \\ \hline
19 & Латвия & 478,00 & 739,31 & 1546,68 \\ \hline
20 & Литва & 947,60 & 585,15 & 617,51 \\ \hline
21 & Люксембург & 615,00 & 6069,44 & 9869,00 \\ \hline
22 & Мальта & 968,00 & 3195,54 & 3301,18 \\ \hline
23 & Молдова & 68,85 & 484,07 & 7030,59 \\ \hline
24 & Нидерланды & 6248,00 & 13677,21 & 2189,05 \\ \hline
25 & Норвегия & 1435,00 & 2316,74 & 1614,45 \\ \hline
26 & Польша & 9722,00 & 10150,00 & 1044,02 \\ \hline
27 & Португалия & 9617,00 & 13904,15 & 1445,79 \\ \hline
28 & Румыния & 6788,90 & 3524,78 & 519,20 \\ \hline
29 & Сан-Марино & 94,00 & 224,76 & 2391,06 \\ \hline
30 & Северная Македония & 294,00 & 386,68 & 1315,23 \\ \hline
31 & Сербия & 871,00 & 2179,16 & 2501,90 \\ \hline
32 & Словения & 1832,00 & 2041,47 & 1114,34 \\ \hline
33 & Турция & 29925,00 & 35209,00 & 1176,57 \\ \hline
34 & Украина & 3973,00 & 1387,00 & 349,11 \\ \hline
35 & Франция & 48395,00 & 44430,27 & 918,08 \\ \hline
36 & Хорватия & 10641,00 & 10899,96 & 1024,34 \\ \hline
37 & Черногория & 1554,00 & 902,90 & 581,02 \\ \hline
38 & Чехия & 3768,00 & 3932,47 & 1043,65 \\ \hline
39 & Швейцария & 4390,00 & 13069,16 & 2977,03 \\ \hline
40 & Швеция & 2990,00 & 6056,00 & 2025,42 \\ \hline
41 & Эстония & 808,00 & 896,39 & 1109,40 \\ \hline
\end{longtable}
\normalsize

Даже при простом просмотре таблицы видно, что значения заметно различаются. В нижней части распределения находятся страны, где на одно прибытие приходится менее 600 USD, а в верхней части выборки есть наблюдения выше 3000 USD. Уже на этом этапе можно предположить наличие выраженной асимметрии распределения и отдельных выбросов.

\subsection{Наглядное представление данных}

Для первичного визуального анализа была построена гистограмма распределения показателя \textit{Receipts per Arrival}. Такой график помогает понять, где сосредоточена основная масса наблюдений и есть ли у распределения длинный правый хвост.

По гистограмме (рис.~\ref{fig:hist}) видно, что большая часть стран сосредоточена в области примерно до 2000 USD на одно прибытие. При этом несколько стран имеют заметно более высокие значения. Из-за этого правая часть распределения оказывается растянутой.

\begin{figure}[H]
    \centering
    \insertfig{\figHistPath}
    \caption{Гистограмма распределения показателя \textit{Receipts per Arrival} по странам Европы за 2021 год}
    \label{fig:hist}
\end{figure}

Уже на этом этапе можно сделать предварительный вывод: рассматриваемый показатель не выглядит симметричным. Для него естественнее предположить модель с положительными значениями и выраженным правым хвостом.

\subsection{Выборочные оценки числовых характеристик данных}

Для описания выборки были рассчитаны основные числовые характеристики: минимум, нижний квартиль, медиана, верхний квартиль, максимум, размах, среднее значение, дисперсия, стандартное отклонение и коэффициент вариации.

Выборочное среднее определяется формулой
\begin{equation}
\overline{x} = \frac{1}{n} \sum_{i=1}^n x_i,
\end{equation}
а выборочная дисперсия - формулой
\begin{equation}
s^2 = \frac{1}{n-1} \sum_{i=1}^n \left(x_i - \overline{x}\right)^2.
\end{equation}

Стандартное отклонение равно
\begin{equation}
s = \sqrt{s^2},
\end{equation}
а коэффициент вариации вычисляется как
\begin{equation}
V = \frac{s}{\overline{x}} \cdot 100\%.
\end{equation}

Полученные значения сведены в таблицу~\ref{tab:desc}.

\begin{table}[H]
\centering
\caption{Основные числовые характеристики выборки}
\label{tab:desc}
\begin{tabular}{|p{7.5cm}|p{4cm}|}
\hline
Характеристика & Значение \\ \hline
Объем выборки, $n$ & 41 \\ \hline
Минимум, USD & 349,11 \\ \hline
Нижний квартиль $Q_1$, USD & 928,57 \\ \hline
Медиана, USD & 1109,40 \\ \hline
Верхний квартиль $Q_3$, USD & 2025,42 \\ \hline
Максимум, USD & 9869,00 \\ \hline
Размах, USD & 9519,90 \\ \hline
Среднее значение, USD & 1797,49 \\ \hline
Стандартное отклонение, USD & 1856,37 \\ \hline
Выборочная дисперсия & 3446106,23 \\ \hline
Коэффициент вариации, \% & 103,28 \\ \hline
\end{tabular}
\end{table}

Среднее значение составило 1797,49 USD. Это означает, что если мысленно усреднить все страны выборки, то на одно международное прибытие пришлось бы около 1797 долларов туристских поступлений. Однако опираться только на среднее в данном случае не очень удобно, поскольку несколько крупных значений заметно смещают его вверх.

Медиана равна 1109,40 USD. Она делит упорядоченную выборку пополам: у половины стран показатель меньше этого уровня, у второй половины - больше. В данной задаче медиана лучше отражает типичный уровень показателя, чем обычное среднее.

Стандартное отклонение составило 1856,37 USD, а коэффициент вариации - 103,28\%. Это очень высокий уровень вариации. Иначе говоря, различия между странами не сводятся к небольшим отклонениям около общего среднего, а действительно являются существенными. По сочетанию более высокого среднего значения по сравнению с медианой и очень большого коэффициента вариации можно сделать вывод о правосторонней асимметрии распределения.

\subsection{Оценка плотности распределения и проверка гипотезы о характере распределения}

Поскольку все наблюдения положительны, а гистограмма имеет правый хвост, в качестве рабочей гипотезы выбрано логнормальное распределение. Такое предположение достаточно типично для экономических и стоимостных показателей: величина не может быть отрицательной, а редкие высокие значения при этом допустимы.

Далее рассмотрим логарифмы исходных значений:
\begin{equation}
y_i = \ln x_i.
\end{equation}

Для них оцениваются параметры нормального распределения:
\begin{equation}
\hat{\mu} = \frac{1}{n} \sum_{i=1}^n y_i,
\end{equation}
\begin{equation}
\hat{\sigma} = \sqrt{\frac{1}{n-1} \sum_{i=1}^n \left(y_i - \overline{y}\right)^2}.
\end{equation}

По выборке получены оценки
\[
\hat{\mu} = 7,197, \qquad \hat{\sigma} = 0,713.
\]

Тогда оценка плотности логнормального распределения записывается в виде
\begin{equation}
f(x) = \frac{1}{x \hat{\sigma} \sqrt{2\pi}}
\exp\left(
-\frac{(\ln x - \hat{\mu})^2}{2\hat{\sigma}^2}
\right), \qquad x>0.
\end{equation}

На рисунке~\ref{fig:lognorm} показана гистограмма с наложенной оценкой плотности логнормального распределения. По виду графика видно, что предложенная модель достаточно хорошо повторяет общую форму эмпирического распределения: максимум находится в области сравнительно небольших значений, а затем плотность постепенно убывает вправо.

\begin{figure}[H]
    \centering
    \insertfig{\figLognormPath}
    \caption{Гистограмма распределения показателя \textit{Receipts per Arrival} с наложенной оценкой плотности логнормального распределения}
    \label{fig:lognorm}
\end{figure}

Для формальной проверки гипотезы применен критерий согласия Пирсона $\chi^2$. Число интервалов принято равным $k=6$. Такой выбор позволяет, с одной стороны, не делать группировку слишком грубой, а с другой - выполнить стандартное условие применимости критерия: ожидаемые частоты не должны быть слишком малыми. Интервалы выбирались по квантилям подогнанного логнормального распределения, поэтому ожидаемая частота в каждом интервале одинакова:
\[
E_i = \frac{n}{k} = \frac{41}{6} \approx 6,83.
\]

Результаты расчета приведены в таблице~\ref{tab:chi}.

\begin{table}[H]
\centering
\caption{Интервалы и частоты для критерия Пирсона}
\label{tab:chi}
\begin{tabular}{|p{5.5cm}|c|c|c|}
\hline
Интервал значений, USD & $O_i$ & $E_i$ & $\dfrac{(O_i-E_i)^2}{E_i}$ \\ \hline
0,00 -- 675,47 & 6 & 6,83 & 0,102 \\ \hline
675,47 -- 985,62 & 8 & 6,83 & 0,199 \\ \hline
985,62 -- 1334,79 & 11 & 6,83 & 2,541 \\ \hline
1334,79 -- 1807,65 & 3 & 6,83 & 2,150 \\ \hline
1807,65 -- 2637,67 & 8 & 6,83 & 0,199 \\ \hline
свыше 2637,67 & 5 & 6,83 & 0,492 \\ \hline
\end{tabular}
\end{table}

Статистика критерия вычисляется по формуле
\begin{equation}
\chi^2_{\text{набл}} = \sum_{i=1}^k \frac{(O_i-E_i)^2}{E_i}.
\end{equation}

В результате получено
\[
\chi^2_{\text{набл}} = 5,683.
\]

Число степеней свободы равно
\[
\nu = k - 1 - m = 6 - 1 - 2 = 3,
\]
где $m=2$ - число оцененных параметров распределения.

При уровне значимости $\alpha = 0{,}05$ критическое значение составляет
\[
\chi^2_{\text{кр}} = 7,815.
\]

Так как
\[
\chi^2_{\text{набл}} = 5,683 < \chi^2_{\text{кр}} = 7,815,
\]
то оснований отвергать гипотезу о логнормальном распределении нет. Дополнительно можно отметить, что $p$-значение равно 0,128, то есть тоже больше уровня значимости 0,05.

Следовательно, для данной выборки логнормальную модель можно считать приемлемой. Это не означает, что распределение описано идеально, однако для учебного статистического анализа такая гипотеза достаточно хорошо согласуется с имеющимися данными.

\subsection{Интервальные оценки числовых характеристик данных}

Сначала построим приближенный 95\%-й доверительный интервал для среднего значения исходного показателя. При объеме выборки $n=41$ можно воспользоваться $t$-распределением:
\begin{equation}
\overline{x} \pm t_{1-\alpha/2;\,n-1} \cdot \frac{s}{\sqrt{n}}.
\end{equation}

Для исследуемой выборки получаем интервал
\[
(1211,55;\ 2383,43) \text{ USD}.
\]

Этот интервал достаточно широк, что хорошо согласуется с уже отмеченным большим разбросом данных.

Так как гипотеза о логнормальном распределении не отвергнута, дополнительно построим интервальные оценки для логарифмированных данных. Для параметра $\mu$ нормального распределения логарифмов используется формула
\begin{equation}
\overline{y} \pm t_{1-\alpha/2;\,n-1} \cdot \frac{s_y}{\sqrt{n}}.
\end{equation}

Получаем интервал
\[
(6,972;\ 7,422).
\]

После обратного преобразования получаем доверительный интервал для медианы логнормального распределения:
\[
\left(e^{6.972};\ e^{7.422}\right)
=
(1065,86;\ 1671,57) \text{ USD}.
\]

Иначе говоря, с доверительной вероятностью 0,95 типичный уровень показателя, если ориентироваться на медиану логнормального распределения, находится примерно между 1066 и 1672 USD.

Для дисперсии логарифмов строится интервал
\begin{equation}
\frac{(n-1)s_y^2}{\chi^2_{1-\alpha/2;\,n-1}}
\le \sigma_y^2 \le
\frac{(n-1)s_y^2}{\chi^2_{\alpha/2;\,n-1}}.
\end{equation}

В численном виде:
\[
(0,342;\ 0,832).
\]

Таким образом, интервальные оценки подтверждают тот же вывод, что и выборочные характеристики: показатель имеет не только асимметричную форму, но и достаточно заметную вариацию.

\subsection{Выводы по основной части}

В основной части работы была сформирована выборка по 41 стране Европы и рассчитан показатель \textit{Receipts per Arrival} за 2021 год. По гистограмме и числовым характеристикам установлено, что распределение имеет выраженную правостороннюю асимметрию. Среднее значение оказалось больше медианы, а коэффициент вариации превысил 100\%.

Гипотеза о логнормальном распределении была проверена критерием Пирсона и не была отвергнута. Это дало возможность использовать логарифмирование данных при построении интервальных оценок. В результате можно сделать вывод, что для рассматриваемой совокупности стран логнормальная модель является разумным приближением.

\section{Дополнительный анализ}

\subsection{Сравнение стран по уровню показателя}

Чтобы анализ не оставался слишком абстрактным, полезно обратиться к конкретным странам. Наибольшие значения показателя в 2021 году наблюдаются у Люксембурга (9869,00 USD), Молдовы (7030,59 USD), Великобритании (6010,89 USD), Мальты (3301,18 USD) и Швейцарии (2977,03 USD).

Минимальные значения, наоборот, относятся к Украине (349,11 USD), Албании (449,67 USD), Румынии (519,20 USD), Черногории (581,02 USD) и Дании (585,94 USD).

Такой результат показывает, что совокупность европейских стран в выборке очень неоднородна. Высокое значение показателя не обязательно означает большое число туристов; иногда оно связано с высокой средней стоимостью поездки, более длительным пребыванием, более дорогой структурой услуг или сравнительно небольшим потоком туристов при ощутимом объеме поступлений. Именно поэтому показатель \textit{Receipts per Arrival} полезно анализировать отдельно от абсолютного числа прибытий.

Для наглядности ниже приведена столбчатая диаграмма десяти стран с наибольшими значениями показателя (рис.~\ref{fig:top10}). Она хорошо показывает, насколько сильно верхняя часть распределения отрывается от основной массы наблюдений.

\begin{figure}[H]
    \centering
    \insertfig{\figTopPath}
    \caption{Десять стран с наибольшим значением \textit{Receipts per Arrival} в 2021 году}
    \label{fig:top10}
\end{figure}

\subsection{Сравнение укрупненных групп стран}

В качестве дополнительного анализа выполним укрупненную группировку стран по подрегионам Европы. Эта группировка является рабочей и используется не для строгой географической классификации, а для того, чтобы сравнить общие тенденции между разными группами стран.

Результаты приведены в таблице~\ref{tab:subregions}.

\begin{table}[H]
\centering
\caption{Сравнение укрупненных групп стран}
\label{tab:subregions}
\small
\begin{tabular}{|p{4.2cm}|>{\centering\arraybackslash}p{0.8cm}|>{\centering\arraybackslash}p{2.0cm}|>{\centering\arraybackslash}p{2.0cm}|>{\centering\arraybackslash}p{1.9cm}|>{\centering\arraybackslash}p{1.9cm}|}
\hline
Группа стран & $n$ & Среднее, USD & Медиана, USD & Минимум, USD & Максимум, USD \\ \hline
Северная Европа & 8 & 1922,10 & 1580,56 & 585,94 & 6010,89 \\ \hline
Западная Европа & 8 & 2792,51 & 2044,46 & 918,08 & 9869,00 \\ \hline
Южная Европа & 13 & 1357,72 & 1096,06 & 449,67 & 3301,18 \\ \hline
Восточная и Юго-Восточная Европа & 12 & 1527,50 & 986,11 & 349,11 & 7030,59 \\ \hline
\end{tabular}
\normalsize
\end{table}

По среднему значению лидирует западноевропейская группа (2792,51 USD). Такой результат во многом объясняется присутствием Люксембурга, Швейцарии, Бельгии и Нидерландов с высокими значениями показателя. На втором месте находится Северная Европа (1922,10 USD), где на итог заметно влияет Великобритания.

Для Южной Европы среднее значение ниже - 1357,72 USD, а для Восточной и Юго-Восточной Европы - 1527,50 USD. При этом по медиане различие между южной и восточной группами не выглядит очень большим: 1096,06 и 986,11 USD соответственно. Это еще раз показывает, что в некоторых группах общую картину существенно меняют отдельные страны с аномально высокими значениями.

\subsection{Сравнение с 2019 годом}

Для дополнительного содержательного анализа те же данные были сопоставлены с 2019 годом. Сравнение выполнялось только по тем странам, для которых показатель удалось рассчитать и за 2019, и за 2021 год. Таких стран оказалось 31. Это важно: при сравнении по времени состав совокупности должен быть одинаковым, иначе результат может исказиться из-за смены набора стран.

Сводные результаты приведены в таблице~\ref{tab:2019_2021}.

\begin{table}[H]
\centering
\caption{Сравнение показателя \textit{Receipts per Arrival} в 2019 и 2021 годах}
\label{tab:2019_2021}
\small
\begin{tabular}{|p{6.8cm}|>{\centering\arraybackslash}p{2.2cm}|>{\centering\arraybackslash}p{2.2cm}|}
\hline
Характеристика & 2019 год & 2021 год \\ \hline
Объем выборки & 31 & 31 \\ \hline
Среднее значение, USD & 1042,47 & 1733,44 \\ \hline
Медиана, USD & 778,60 & 1084,04 \\ \hline
Стандартное отклонение, USD & 988,29 & 1934,30 \\ \hline
Коэффициент вариации, \% & 94,80 & 111,59 \\ \hline
\end{tabular}
\normalsize
\end{table}

По общей картине видно, что в 2021 году показатель вырос. Среднее значение увеличилось с 1042,47 до 1733,44 USD, то есть примерно на 66,3\%. Медиана выросла с 778,60 до 1084,04 USD, то есть примерно на 39,2\%.

Примечательно, что во всей сопоставимой выборке не нашлось ни одной страны, где значение показателя в 2021 году оказалось бы ниже уровня 2019 года. Это можно интерпретировать следующим образом: после пандемийного шока международные поездки восстанавливались неравномерно, и в среднем на одно прибытие в 2021 году приходилось больше поступлений, чем в допандемийный период. Возможное объяснение состоит в том, что раньше вернулся не массовый, а более платежеспособный спрос. Кроме того, на показатель могли влиять инфляционные и ценовые факторы.

Для визуального сопровождения этого сравнения удобно использовать диаграмму изменения показателя по странам (рис.~\ref{fig:compare}).

\begin{figure}[H]
    \centering
    \insertfig{\figComparePath}
    \caption{Изменение показателя \textit{Receipts per Arrival} по странам, для которых есть данные за 2019 и 2021 годы}
    \label{fig:compare}
\end{figure}

\subsection{Диаграмма размаха и дополнительные числовые показатели}

Для выявления выбросов и уточнения структуры распределения построена диаграмма размаха. Такой график особенно удобен в тех случаях, когда в выборке есть несколько явно выделяющихся наблюдений.

\begin{figure}[H]
    \centering
    \insertfig[0.75\textwidth]{\figBoxPath}
    \caption{Диаграмма размаха распределения показателя \textit{Receipts per Arrival}}
    \label{fig:boxplot}
\end{figure}

Для рассматриваемой выборки нижний квартиль равен 928,57 USD, верхний квартиль - 2025,42 USD. Межквартильный размах составляет
\[
IQR = Q_3 - Q_1 = 1096,85 \text{ USD}.
\]

Верхняя граница по правилу $1{,}5 \cdot IQR$ равна
\[
Q_3 + 1{,}5IQR = 3670,69 \text{ USD}.
\]

Следовательно, к выбросам на диаграмме размаха относятся страны, значения которых превышают 3671 USD. В данной выборке это Люксембург, Молдова и Великобритания. Этот результат хорошо согласуется и с таблицей исходных данных, и с гистограммой распределения.

Дополнительно были рассчитаны еще несколько характеристик, которые помогают описать форму распределения более подробно (таблица~\ref{tab:robust}).

\begin{table}[H]
\centering
\caption{Дополнительные числовые характеристики распределения}
\label{tab:robust}
\begin{tabular}{|p{8cm}|p{3.5cm}|}
\hline
Характеристика & Значение \\ \hline
Геометрическое среднее, USD & 1334,79 \\ \hline
10\%-ное усеченное среднее, USD & 1381,41 \\ \hline
Коэффициент асимметрии Пирсона & 1,112 \\ \hline
Коэффициент асимметрии (моментный) & 2,988 \\ \hline
Коэффициент эксцесса & 9,829 \\ \hline
\end{tabular}
\end{table}

Геометрическое среднее (1334,79 USD) и 10\%-ное усеченное среднее (1381,41 USD) заметно ниже обычного арифметического среднего. Это ожидаемо: обе эти характеристики слабее реагируют на крупные выбросы и поэтому лучше передают типичный уровень для основной массы стран.

Коэффициент асимметрии Пирсона равен 1,112. Так как он положителен и больше единицы, можно говорить о выраженной правосторонней асимметрии. Моментный коэффициент асимметрии (2,988) и коэффициент эксцесса (9,829) подтверждают тот же вывод: распределение не только асимметрично, но и имеет более тяжелый правый хвост, чем нормальное.

\subsection{Выводы по дополнительной части}

Дополнительный анализ сделал общую картину более предметной. Во-первых, были выделены конкретные страны, задающие верхнюю и нижнюю части распределения. Во-вторых, сравнение укрупненных групп показало, что наиболее высокие значения характерны для западноевропейской группы, тогда как южная и восточная группы в среднем выглядят более умеренными.

Наконец, сопоставление с 2019 годом показало общий рост показателя в 2021 году по всей сопоставимой выборке. Это делает интерпретацию результатов более содержательной: различия между странами объясняются не только географическим положением, но и особенностями восстановления туристского рынка после кризисного периода.

\section*{Заключение}
\addcontentsline{toc}{section}{Заключение}

В ходе работы был проведен статистический анализ показателя \textit{Receipts per Arrival} по 41 стране Европы за 2021 год на основе данных UN Tourism. Показатель был рассчитан самостоятельно по двум исходным таблицам - по прибытиям и по поступлениям от въездного туризма.

Расчет основных числовых характеристик показал, что выборка неоднородна: среднее значение составило 1797,49 USD, медиана - 1109,40 USD, коэффициент вариации - 103,28\%. Это свидетельствует о большом разбросе и существенных различиях между странами. Гистограмма и диаграмма размаха показали правостороннюю асимметрию и наличие нескольких выраженных выбросов.

Гипотеза о логнормальном распределении была проверена критерием Пирсона и не была отвергнута при уровне значимости 0,05. Это позволило дополнительно построить интервальные оценки на логарифмированной шкале и получить более устойчивую интерпретацию типичного уровня показателя.

Дополнительный раздел анализа показал, что внутри Европы различия по данному показателю весьма существенны. В верхней части распределения находятся страны, где на одно прибытие приходится особенно большой объем поступлений, тогда как в нижней части выборки имеются страны с заметно более низким средним расходом на одно прибытие. Сравнение с 2019 годом показало общий рост показателя по сопоставимой совокупности стран.

В целом можно сделать вывод, что в 2021 году европейские страны очень сильно различались по величине туристских поступлений на одно международное прибытие. Поэтому для содержательной оценки въездного туризма важно учитывать не только число прибытий, но и средние поступления на одну поездку.

\section*{Приложение}
\addcontentsline{toc}{section}{Приложение}

\begin{enumerate}[leftmargin=1.25cm]
    \item SPoI lab-1 Github. URL: \url{https://github.com/Echways/SPoI}.
\end{enumerate}

\end{document}